\pdfbookmark[2]{Zu VI und VII.}{nexus-remarks-vi-vii}
\begin{center}\large Zu VI und VII.\end{center}

Wie der Verfasser die Reihe $\Sigma\left(\frac{-D}n\right)\frac1n$ summiert hat, geht aus einigen besonderen Beispielen hervor, die sich auf einzelnen Blättern finden.

Ist $D\equiv3\pmod4$, so folgt aus dem Fundamentaltheorem der Theorie der quadratischen Reste unter Verwendung der Reihe
\[
\operatorname{cotang}u=\frac1u+\frac1{u-\pi}+\frac1{u+\pi}+\frac1{u-2\pi}+\frac1{u+2\pi}+\ldots
\]
dass
\[
\Sigma\left(\frac{-D}n\right)\frac1n=\Sigma\left(\frac nD\right)\frac1n=\frac\pi{2D}\Sigma\left(\frac\nu D\right)\operatorname{cotang}\frac{\nu\pi}{2D}
\]
ist, wobei $\nu$ alle zu $2D$ teilerfremden Zahlen durchläuft, die kleiner als $D$ sind. Setzt man
\[
\surd-1=i,\quad\cos\frac{2\pi}D+i\sin\frac{2\pi}D=r
\]
und bezeichnet mit $\mu$ alle zu $D$ teilerfremden Zahlen, die nicht größer als $D$ sind, so lässt sich die vorstehende Summe leicht umformen zu
\[
\Sigma\left(\frac{-D}n\right)\frac1n=\frac{\pi i}{4D}\left(\frac2D\right)\Sigma\left(\frac\mu D\right)\frac{r^\mu-1}{r^\mu+1}
\]
Wendet man nun die für jede Wurzel $\omega$ der Gleichung $\omega^D=1$ gültige Formel
\SourceLeaf{308}
\[
\frac{\omega-1}{\omega+1}=\Sigma(-1)^{\alpha-1}\omega^\alpha
\]
an, in der $\alpha$ die Zahlen $1,2,3\ldots(D-1)$ durchlaufen muss, so erhält man durch Umkehrung der Summationsreihenfolge
\[
\Sigma\left(\frac{-D}n\right)\frac1n=\frac{\pi i}{4D}\left(\frac2D\right)\Sigma(-1)^{\alpha-1}\Sigma\left(\frac\mu D\right)r^{\alpha\mu}
\]
\Needspace{9\baselineskip}
Die Summation über $\mu$ lässt sich bekanntlich mithilfe der in der Abhandlung \emph{Summatio quarundam serierum singularium} bewiesenen Sätze ausführen. Beschränkt man sich auf den Fall, in dem $D$ durch kein Quadrat teilbar ist, so findet man allgemein
\begingroup\postdisplaypenalty=10000
\[
\Sigma\left(\frac\mu D\right)r^{\alpha\mu}=\left(\frac\alpha D\right)i^{\left(\frac{D-1}2\right)^2}\surd D
\]
\endgroup
wobei $\left(\frac\alpha D\right)=0$ zu setzen ist, wenn $\alpha$ nicht zu $D$ teilerfremd ist. Im Fall $D\equiv3\pmod4$ erhält man daher
\begingroup\postdisplaypenalty=10000\interlinepenalty=10000
\[
\Sigma\left(\frac{-D}n\right)\frac1n=\frac\pi{4\surd D}\left(\frac2D\right)\Sigma(-1)^\alpha\left(\frac\alpha D\right)=\frac\pi{2\surd D}\Sigma\left(\frac{\alpha^{\prime}}D\right)
\]
wobei $\alpha^{\prime}$ alle zu $D$ teilerfremden Zahlen durchläuft, die kleiner als $\frac12D$ sind. Da schließlich $\epsilon=2$ ist, wird die Klassenanzahl
\[
k=\Sigma\left(\frac{\alpha^{\prime}}D\right)
\]
\endgroup

Ist dagegen $D\equiv1\pmod4$, so erhält man unter Verwendung der Reihe
\[
\operatorname{cosec}u=\frac1u-\frac1{u-\pi}-\frac1{u+\pi}+\frac1{u-2\pi}+\frac1{u+2\pi}-\ldots
\]
auf ähnliche Weise
\[
\begin{aligned}
\Sigma\left(\frac{-D}n\right)\frac1n&=\Sigma(-1)^{\frac{n-1}2}\left(\frac nD\right)\frac1n
=\frac\pi{2D}\Sigma(-1)^{\frac{\nu-1}2}\left(\frac\nu D\right)\operatorname{cosec}\frac{\nu\pi}{2D}\\
&=\frac\pi{2D}\Sigma\left(\frac\mu D\right)\frac{r^\mu}{r^{2\mu}+1}
\end{aligned}
\]
wobei die Buchstaben $\nu$ und $\mu$ ihre frühere Bedeutung haben. Schließt man den offensichtlichen Fall $D=1$ aus und wendet die für jede Wurzel $\omega$ der Gleichung $\omega^D=1$ (außer $\omega=1$) gültige Formel
\[
\frac\omega{\omega\omega+1}=1+\Sigma\omega^{4\alpha^{\prime\prime}}+\Sigma\omega^{D-4\alpha^{\prime\prime}}
\]
an, in der $\alpha^{\prime\prime}$ die Zahlen $1,2,3\ldots\frac14(D-1)$ durchlaufen muss, so ergibt sich, wieder unter der Einschränkung, dass $D$ durch kein Quadrat teilbar ist,
\[
\Sigma\left(\frac{-D}n\right)\frac1n=\frac\pi{\surd D}\Sigma\left(\frac{\alpha^{\prime\prime}}D\right)
\]
und daraus, da $\epsilon=2$ ist,
\[
k=2\Sigma\left(\frac{\alpha^{\prime\prime}}D\right)
\]

\SourceLeaf{309}
Die Fälle, in denen $D$ gerade ist, ließen sich ganz ähnlich behandeln. —

\Needspace{19\baselineskip}
Zur Bestimmung der Klassenanzahl für \emph{positive} Determinanten $D$ finden sich außer der im Text mitgeteilten Schlussformel nur einzelne geometrische Figuren. Sie stellen Hyperbelsektoren endlicher Ausdehnung dar; daneben stehen Ungleichungen, die die Punkte, deren Koordinaten die Variablen der quadratischen Formen sind, auf das Innere eines solchen Hyperbelsektors beschränken. Diese Hyperbelsektoren treten an die Stelle der Ellipsen, die den quadratischen Formen mit negativen Determinanten entsprechen. Durch die Bestimmung ihrer Flächeninhalte ergibt sich wiederum die mittlere Darstellungsanzahl, sofern nur Darstellungen zugelassen werden, deren Variablen die erwähnten Ungleichungen erfüllen. Andererseits dienen diese Ungleichungen dazu, aus den unendlich vielen Darstellungen einer Zahl $m$, die alle zu ein und derselben Wurzel $n$ der Kongruenz $nn-D\equiv0\pmod{\frac m{\mu\mu}}$ gehören und sämtlichen Lösungen der Gleichung $tt-Duu=1$ entsprechen (vgl. Disqq. Arithm., Art. 205), eine einzige herauszugreifen und alle anderen auszuschließen. Die Anzahl aller zugelassenen Darstellungen der Zahl $m$ durch die Gesamtheit aller nicht eigentlich äquivalenten, eigentlich primitiven Formen ist dann gleich dem Wert der Funktion $f(m)$, in der nur $-D$ durch $D$ zu ersetzen ist. Aus ihren Eigenschaften ergibt sich wie zuvor bei negativen Determinanten ein zweiter Ausdruck für die mittlere Darstellungsanzahl; der Vergleich mit dem zuvor geometrisch abgeleiteten Wert führt dann unmittelbar zur Bestimmung der Klassenanzahl.
\BookRule{35mm}
